% Independent derivation for the root's actual-adelic/coinvariant bridge.
% Algebraic quotients only; no assertion about an unproved topological quotient.
\section{The actual adelic source and its retained prime boundary}
\label{sec:independent-adelic-coinvariant-bridge}

The adelic restriction operation and its original measures are those of
Alain Connes, \emph{Trace formula in noncommutative geometry and the zeros
of the Riemann zeta function},
\href{https://arxiv.org/abs/math/9811068v1}{Section III (6)--(19)},
and Alain Connes, Caterina Consani and Matilde Marcolli,
\href{https://arxiv.org/abs/math/0703392v1}{\emph{The Weil proof and the
geometry of the adeles class space}}, as specified in SWQ2--14.
The homological complex used below is the standard Koszul complex,
named for Jean-Louis Koszul. The standard regular-sequence mechanism is
recorded by \emph{The Stacks Project Authors},
\href{https://stacks.math.columbia.edu/tag/062F}{Tag 062F}, with
\href{https://github.com/stacks/stacks-project/blob/a04446e57ec1fbc252a871afcec7752fb2807b14/more-algebra.tex#L7264}{original
author TeX, lines 7264--7298, commit
\texttt{a04446e57ec1fbc252a871afcec7752fb2807b14}}.
The contribution here is the exact application, connecting representatives,
and receiving maps for these specified programme modules; the necessary
Koszul exactness is proved below as well. All tensor products and
quotients below are algebraic unless a different meaning is explicitly
stated.

The precise public programme inputs are
\href{https://github.com/KokunoYumeto/zeta-function-research-reader/blob/acd7016b8232a56f31ae3e69fb800571e8789080/workbenches/splitzero-tandem/continuations/20260923-retained-source-readers/original-zeta-full-boundary/CC_WEIL_BOUNDARY_READER.tex#L927}{SWQ2--14},
\href{https://github.com/KokunoYumeto/zeta-function-research-reader/blob/acd7016b8232a56f31ae3e69fb800571e8789080/workbenches/splitzero-tandem/continuations/20260923-retained-source-readers/original-zeta-full-support/MASS_THETA_PRIME_RECEIVER.tex#L12}{MRT1--20}, and
\href{https://github.com/KokunoYumeto/zeta-function-research-reader/blob/acd7016b8232a56f31ae3e69fb800571e8789080/workbenches/splitzero-tandem/continuations/20260923-retained-source-readers/endpoint-resonance/ENDPOINT_RESONANCE_AND_PRIME_CLASS.tex#L430}{ER27--30}.
These pinned proof targets were checked against the local source versions.

\subsection{The finite adelic basis and the exact original source}

Put
\begin{equation}\tag{ABR1}
 G=\mathbb Q_{>0}^{\times},\quad R=\mathbb C[G],\quad
 I=\ker(\varepsilon:R\longrightarrow\mathbb C),\quad
 H=\mathcal S_{\rm even}(\mathbb R),\quad
 K=\widehat{\mathbb Z}^{\times},\quad c_a=1_{a\widehat{\mathbb Z}}.
\end{equation}
The basis element $t_a\in R$ satisfies $t_at_b=t_{ab}$ and
$\varepsilon(t_a)=1$. The finite additive measure has
$\operatorname{vol}(\mathbb Z_p)=1$, and the real additive measure
is Lebesgue measure. Thus
\begin{equation}\tag{ABR2}
 \operatorname{vol}(a\widehat{\mathbb Z})
  =\prod_p p^{-v_p(a)}=a^{-1}\qquad(a>0\text{ rational}).
\end{equation}
This equality uses the original positive real value of $a$ and its
prime factorization; the finite and real factors will both be retained.

The family $(c_a)_{a\in G}$ is a vector-space basis of
$\mathcal S(\mathbb A_f)^K$. To prove spanning, first observe that a
compactly supported locally constant function on $\mathbb Q_p$,
invariant under $\mathbb Z_p^\times$, is constant on valuation shells.
Choose integers $A\le B$ such that its support is in $p^A\mathbb Z_p$
and it is constant on cosets of $p^B\mathbb Z_p$. If its values on
the shells $A\le n<B$ and on the last ball $p^B\mathbb Z_p$ are
$d_A,\ldots,d_B$, the function is exactly
\[
 d_A1_{p^A\mathbb Z_p}
 +\sum_{n=A+1}^{B}(d_n-d_{n-1})1_{p^n\mathbb Z_p}.
\]
The original Bruhat space on $\mathbb A_f$ consists of finite sums
of restricted tensor products with $1_{\mathbb Z_p}$ outside a finite
set of primes. Compact-unit averaging of a finite-place test is a
finite linear operation: its support is contained in finitely many
bounded local balls and its local-constancy subgroups have finite
index there. Averaging therefore reduces to a finite quotient. Apply
the preceding local formula at the finitely many relevant primes and
expand the finite tensor products. Their balls are exactly $c_a$.
For an invariant test averaging leaves the test itself fixed, which
proves spanning.

For independence, take distinct $a_1,\ldots,a_m\in G$ and partially
order them by $a_i\preceq a_j$ when $a_j/a_i$ is an integer. This is
a partial order because both positive ratios can be integers only
when they are one. Order this finite set by a linear extension.
Evaluate its characteristic functions at the finite rational ideles
$x_f=a_j$. Then
\[
 c_{a_i}(a_j)=1\quad\Longleftrightarrow\quad a_i\preceq a_j.
\]
Indeed $\mathbb Q\cap\widehat{\mathbb Z}=\mathbb Z$.
The resulting matrix is lower triangular with diagonal one. Its
determinant is one, proving independence. This proof also works for
coefficients in any complex vector space, because back substitution
uses only this finite triangular matrix.

Define
\begin{equation}\tag{ABR3}
 \Theta:R\otimes H\longrightarrow
           \mathcal S(\mathbb A_{\mathbb Q})^{K,\mathrm{even}},
 \qquad
 \Theta(t_a\otimes h)(x_\infty,x_f)
       =h(x_\infty/a)c_a(x_f).
\end{equation}
Here evenness refers to the original real variable $x_\infty$.
Every target is a finite sum $\sum_a g_a(x_\infty)c_a(x_f)$ with
$g_a\in H$. Set $h_a(x)=g_a(ax)$ to obtain its preimage in ABR3.
Independence follows from the vector-valued independence of the
$c_a$ and the injectivity of each real dilation. Thus $\Theta$ is
an isomorphism of vector spaces, with the displayed exact inverse.
No change in the original real coordinate has been made in the target.

For $r\in G$ define full rational adelic dilation by
$\alpha_r\xi(x)=\xi(r^{-1}x)$. The finite argument gives
$c_a(r^{-1}x_f)=c_{ra}(x_f)$, while the real argument is
$h(x_\infty/(ra))$. Consequently
\begin{equation}\tag{ABR4}
 \alpha_r\Theta(t_a\otimes h)=\Theta(t_{ra}\otimes h).
\end{equation}
This is the left $R$-action. It is full diagonal rational dilation;
it is not dilation of the real idele-class coordinate alone.

Let $m(h)=(h(0),\int_{\mathbb R}h(x)\,dx)$. Evaluation and
integration of ABR3 give the two exact moments
\begin{equation}\tag{ABR5}
 \Theta(t_a\otimes h)(0)=h(0),\qquad
 \int_{\mathbb A}\Theta(t_a\otimes h)
     =a\left(\int_{\mathbb R}h\right)
                     \prod_p p^{-v_p(a)}=\int_{\mathbb R}h.
\end{equation}
Thus $(\varepsilon_0,\varepsilon_1)\Theta=\varepsilon\otimes m$.
In particular
\begin{equation}\tag{ABR6}
 M_0=\ker(\varepsilon\otimes m),\qquad H_{00}=\ker m,
 \qquad
 \Theta(M_0)=\mathcal S(\mathbb A)_0^{K,\mathrm{even}}.
\end{equation}
This is equality of the actual spaces with their two original
endpoint conditions, not a comparison only after completion.

\subsection{The complete algebraic coinvariant and its boundary}

The two even Schwartz functions
\[
 h_0(x)=(1-2\pi x^2)e^{-\pi x^2},\qquad
 h_1(x)=2\pi x^2e^{-\pi x^2}
\]
have moments $(1,0)$ and $(0,1)$. Indeed their evaluations are
immediate and the real Gaussian integrals are
$\int e^{-\pi x^2}=1$ and $\int x^2e^{-\pi x^2}=1/(2\pi)$.
The section $\sigma(u_0,u_1)=u_0h_0+u_1h_1$ therefore splits $m$.
It gives
\begin{equation}\tag{ABR7}
 M_0=(R\otimes H_{00})\oplus(I\otimes\sigma(\mathbb C^2)),
 \qquad
 Q_0=M_0/IM_0\cong H_{00}\oplus(I/I^2)\otimes\mathbb C^2.
\end{equation}
Indeed multiplication by $I$ has exactly the two components
$I\otimes H_{00}$ and $I^2\otimes\sigma(\mathbb C^2)$.

For clarity the resulting boundary and inverse can be written
without leaving a dependence on the chosen moment section. Put
\begin{equation}\tag{ABR8}
 W=\bigoplus_{p\ \mathrm{prime}}\mathbb C\ell_p,
 \qquad \mathcal B_{\mathrm{pr}}=W\otimes\mathbb C^2.
\end{equation}
The map $[t_a-1]\mapsto\sum_pv_p(a)\ell_p$ is an isomorphism
$I/I^2\to W$. For its proof retain the identity
$t_{ab}-1=(t_a-1)+(t_b-1)+(t_a-1)(t_b-1)$.
Modulo $I^2$ it gives additivity, and inversion gives the negative
of a class. Prime factorization expresses each class as the stated
sum. For independence the functionals
$D_p(\sum_ac_at_a)=\sum_ac_av_p(a)$ satisfy
$D_p(rs)=\varepsilon(r)D_p(s)+\varepsilon(s)D_p(r)$, and hence
vanish on $I^2$. They take value one on $[t_p-1]$ and zero on
the other displayed prime generators.

The canonical coordinates of $c=[\sum_at_a\otimes h_a]\in Q_0$
are
\begin{equation}\tag{ABR9}
 \Phi(c)=\sum_a h_a\in H_{00},\qquad
 \beta(c)=\sum_p\ell_p\otimes\sum_a v_p(a)m(h_a).
\end{equation}
Both maps kill $IM_0$. For $\beta$, apply the preceding derivation
identity to a product $r\sum_at_a\otimes h_a$, $r\in I$;
its value is $\sum_pD_p(r)\ell_p\otimes\sum_am(h_a)=0$ by
the defining moment condition. The inverse of ABR9 sends
\begin{equation}\tag{ABR10}
 (h,\sum_p\ell_p\otimes u_p)
 \longmapsto [t_1\otimes h]
             +\sum_p[(t_p-1)\otimes\sigma(u_p)].
\end{equation}
Only finitely many $u_p$ are nonzero. Changing $\sigma$ changes
the latter expression by an element of $I\otimes H_{00}\subset IM_0$.
To check the inverse on an arbitrary representative, subtract
$t_1\otimes\sum_a h_a$ and write the remainder as
$\sum_a(t_a-1)\otimes h_a$. Replacing each $h_a$ there by
$\sigma m(h_a)$ changes its class by $IM_0$, and ABR8 then
gives exactly ABR10. Thus the short exact sequence
\begin{equation}\tag{ABR11}
 0\longrightarrow\mathcal B_{\mathrm{pr}}
   \xrightarrow{\ j\ }Q_0\xrightarrow{\Phi}H_{00}
   \longrightarrow0
\end{equation}
is canonically split by $h\mapsto[t_1\otimes h]$, and $j$ is
the section-independent second term of ABR10.

\subsection{Actual periodization, injectivity, and exact inverse}

For the original restriction map
$E\xi(u,y)=\sum_{q\in\mathbb Q^\times}\xi(qy,qu)$,
$u\in K$, $y>0$, the finite condition in ABR3 is
$q/a\in\widehat{\mathbb Z}\cap\mathbb Q=\mathbb Z$.
Hence $q=an$ for $n\in\mathbb Z\setminus\{0\}$. It follows that
\begin{equation}\tag{ABR12}
 E\Theta(t_a\otimes h)(u,y)
    =\sum_{n\in\mathbb Z\setminus\{0\}}h(ny)
    =2\sum_{n\ge1}h(ny)=:\mathcal P h(y).
\end{equation}
Evenness of $h$ supplies exactly the last factor two. The real
and finite scales $a$ have both participated in the equality;
no finite restriction has been dropped. All sums are locally
absolutely convergent with every real derivative by Schwartz decay.
For $m=\sum_at_a\otimes h_a\in M_0$, ABR12 is
$E\Theta(m)=\mathcal P(\sum_a h_a)$.

The Fourier convention
$\widehat h(\xi)=\int_{\mathbb R}h(x)e^{-2\pi ix\xi}\,dx$
gives the exact Poisson identity with both endpoints
\begin{equation}\tag{ABR13}
 \mathcal P h(y)
 =y^{-1}\mathcal P\widehat h(1/y)
     +y^{-1}\int_{\mathbb R}h(x)\,dx-h(0).
\end{equation}
For $h\in H_{00}$ both displayed endpoint terms are zero, and
$\widehat h$ also belongs to $H_{00}$ because
$\widehat h(0)=\int h=0$ and $\int\widehat h=h(0)=0$.
At infinity the sums of every Euler derivative of $\mathcal P h$
decay faster than every power: choose an arbitrarily large Schwartz
power $A>1$ and bound their sums by a constant times
$y^{-A}\sum n^{-A}$. ABR13 proves the corresponding estimates
at zero. Therefore $\mathcal P h$ belongs to the original radial
strong Schwartz test space. All powers of $\log y$ are also
controlled by these bounds, so its Mellin transform is entire.

For $\Re s>1$, absolute Fubini gives the original identity
\begin{equation}\tag{ABR14}
 \mathcal M(\mathcal P h)(s)
   =2\zeta(s)\mathcal M_+h(s),\qquad
 \mathcal M_+h(s)=\int_0^\infty h(x)x^{s-1}\,dx.
\end{equation}
Indeed the absolute sum of integrals is
$2\zeta(\Re s)\int_0^\infty|h(x)|x^{\Re s-1}\,dx<\infty$.
The original Euler product proves that $\zeta(s)\ne0$ in this
half-plane: the logarithms of its factors converge absolutely
because $\sum_p\sum_{k\ge1}p^{-k\Re s}/k<\infty$.
Thus division by the actual $2\zeta(s)$ is valid on the exact
domain used below. There is no completed replacement of $\zeta$.

For $\sigma>1$, $g_\sigma(v)=e^{\sigma v}h(e^v)$ is Schwartz
on $\mathbb R$. Every derivative is a finite sum of
$e^{(\sigma+j)v}h^{(j)}(e^v)$. These decay exponentially at
$-\infty$ and faster than any prescribed exponential at
$+\infty$, using the original real Schwartz estimates.
Consequently $\mathcal M_+h(\sigma+it)$ decays faster than
every power of $t$. Fourier inversion in $v$ gives
\begin{equation}\tag{ABR15}
 h(x)=\frac1{4\pi i}
   \int_{\sigma-i\infty}^{\sigma+i\infty}
      \frac{\mathcal M(\mathcal P h)(s)}{\zeta(s)}x^{-s}\,ds,
       \qquad x>0,\quad\sigma>1.
\end{equation}
Every constant is explicit: Mellin inversion has $1/(2\pi i)$,
and ABR14 contributes the additional $1/2$. To justify inversion
directly, multiply the Fourier integrand by $e^{-\epsilon t^2}$.
Its inverse is convolution with
$(4\pi\epsilon)^{-1/2}e^{-v^2/(4\epsilon)}$, which tends to
$g_\sigma$ as $\epsilon\downarrow0$. The proved rapid vertical
decay permits dominated convergence on the Fourier side. This
proves ABR15, including absolute convergence. Evenness recovers
$h(x)$ for $x<0$ and $h(0)=0$ supplies the remaining point.
In particular $\mathcal P$ is injective on $H_{00}$.

There is also an exact arithmetic inverse in the original variable:
\begin{equation}\tag{ABR15a}
 h(x)=\frac12\sum_{n\ge1}\mu(n)\mathcal P h(nx),\qquad x>0.
\end{equation}
For fixed $x>0$ and any $A>1$, Schwartz decay bounds the absolute
double sum after inserting ABR12 by a constant depending on $x,A$
times $\sum_{n,m\ge1}(nm)^{-A}=\zeta(A)^2$.
It may therefore be regrouped by $k=nm$. Its coefficient of $h(kx)$
is $\sum_{n\mid k}\mu(n)$, which is one for $k=1$ and zero for
$k>1$: prime factorization gives the finite product
$\prod_{p\mid k}(1-1)$ in the latter case. This proves ABR15a
with its factor $1/2$ and also proves convergence locally with
every Euler derivative. It is a second exact inverse, not a
truncation or a substitute for ABR14--15's original zeta factor.

Let $V=E\mathcal S(\mathbb A)_0$ be the actual restriction image,
and let $V^K$ be its compact-unit invariant part. Then
\begin{equation}\tag{ABR16}
 V^K=E\mathcal S(\mathbb A)_0^{K,\mathrm{even}}
      =\mathcal P(H_{00}).
\end{equation}
For the nontrivial inclusion, take $F=E\xi\in V^K$. Average
$\xi(x_\infty,vx_f)$ over $v\in K$. This is a finite average
in the original Bruhat space, retains both moments, and has
restriction equal to $F$. Next project to real-even functions
by averaging $x_\infty$ with $-x_\infty$. Under $E$, the real
reflection equals the finite-unit translation $u\mapsto-u$,
as seen by reindexing $q\mapsto-q$ in its absolutely convergent
sum. Since $F$ is $K$-invariant this projection still restricts
to $F$. ABR6 now supplies a preimage in $M_0$, whose sum of
coefficients is in $H_{00}$. This proves ABR16 without replacing
$V$ by its closure or by a topological completion.

Full rational dilation is trivial after periodization, also
directly by reindexing $q/r$ in the rational sum. Therefore
ABR12 induces a surjection on the algebraic coinvariant, with
the exact kernel
\begin{equation}\tag{ABR17}
 0\longrightarrow\mathcal B_{\mathrm{pr}}
    \xrightarrow{\ j\ }Q_0
    \xrightarrow{\ \rho=\mathcal P\Phi\ }V^K
    \longrightarrow0.
\end{equation}
Injectivity of $\mathcal P$ proves that this kernel is exactly
the whole prime boundary, rather than merely containing it.
A faithful joint receiver and its inverse are
\begin{equation}\tag{ABR18}
 \Omega(c)=(\rho(c),\beta(c)),\qquad
 \Omega:Q_0\xrightarrow{\ \sim\ }V^K\oplus\mathcal B_{\mathrm{pr}},
 \qquad
 \Omega^{-1}(F,b)=[t_1\otimes\mathcal P^{-1}F]+j(b).
\end{equation}
The inverse $\mathcal P^{-1}$ is exactly ABR15.

Keep Connes's original seminorm, without changing its measure,
factor or weight:
\begin{equation}\tag{ABR19}
 \|c\|_\delta^2
  =\int_0^\infty|y^{1/2}\rho(c)(y)|^2
            (1+\log^2y)^{\delta/2}\frac{dy}{y},\qquad\delta>1.
\end{equation}
It is finite by the strong test estimates. Its weight is strictly
positive. If the integral is zero, continuity gives $\rho(c)=0$
everywhere, and ABR17 gives $c\in j(\mathcal B_{\mathrm{pr}})$.
The reverse inclusion follows from ABR17 as well. Thus ABR19
has precisely this algebraic nullspace. This identifies the
lost coordinates but assigns no new norm to them and claims
no equivalence with an unproved topological quotient.

\subsection{Positive real ideles and the two endpoint characters}

For $b>0$, now with no rationality assumption, define
$R_bh(x)=h(x/b)$ and let $\mathscr R_b=\operatorname{id}_R\otimes R_b$
be the operator on $R\otimes H$.
Its original moments obey
\begin{equation}\tag{ABR19a}
 m(R_bh)=D_bm(h),\qquad
 D_b=\begin{pmatrix}1&0\\0&b\end{pmatrix}.
\end{equation}
Evaluation at zero is unchanged, and the real substitution $x=bu$
multiplies the additive integral by $b$. Thus $\mathscr R_b$
preserves $M_0$, commutes with the full rational action, and
descends to $Q_0$. Under $\Theta$ it is precisely the real-idele
action $\xi(x_\infty,x_f)\mapsto\xi(x_\infty/b,x_f)$; the finite
support $a\widehat{\mathbb Z}$ has not changed.

In the exact coordinates ABR9, this action is
\begin{equation}\tag{ABR19b}
 (\Phi,\beta)\mathscr R_b c
   =\left(R_b\Phi c,
                (\operatorname{id}_W\otimes D_b)\beta c\right).
\end{equation}
This follows directly from the finite sums in ABR9. On the boundary
inverse ABR10 it is also independent of the moment section:
$R_b\sigma(u)-\sigma(D_bu)\in H_{00}$, so its product with
$t_p-1$ is in $I\otimes H_{00}\subset IM_0$. Consequently
$\Omega$ is equivariant for real dilation on its $V^K$ component
and the two endpoint characters $1,b$ on its boundary component.
These are the actual Connes endpoint characters
$\mathbb C\oplus\mathbb C(1)$, with their displayed action;
no purity theorem is inferred from this character calculation.

Both periodization conventions and their exact comparison are
\begin{equation}\tag{ABR19c}
 \begin{split}
 E_{2007}(R_bh)(y)&=(\mathcal P h)(y/b),\\
 E_{1998}(h)(y)&=y^{1/2}\mathcal P h(y),\\
 E_{1998}(R_bh)(y)&=b^{1/2}E_{1998}(h)(y/b).
 \end{split}
\end{equation}
The first equality follows term by term from ABR12. The last one
retains $y^{1/2}=b^{1/2}(y/b)^{1/2}$ and therefore its exact
factor $b^{1/2}$. Multiplication by the nowhere-zero function
$y^{1/2}$ is explicitly inverted by multiplication by $y^{-1/2}$
on the original strong test spaces; it is not omitted in the
seminorm. In particular that seminorm transforms exactly as
\begin{equation}\tag{ABR19d}
 \|\mathscr R_b c\|_\delta^2
   =b\int_0^\infty|E_{1998}(\Phi c)(y)|^2
       \{1+(\log y+\log b)^2\}^{\delta/2}\frac{dy}{y}.
\end{equation}
This follows by the positive substitution in ABR19 and does not
claim that the original weighted seminorm is dilation-invariant.

\subsection{Every higher-degree connecting representative}

Identify $R$ with the Laurent polynomial algebra
$\mathbb C[t_p^{\pm1}:p\text{ prime}]$, whose polynomials involve
only finitely many variables. Let $x_p=t_p-1$ and use the ordered
prime basis $\ell_p$ of $W$. The standard algebraic Koszul complex is
\begin{equation}\tag{ABR20}
 K_n(R)=R\otimes\Lambda^nW,\qquad
 d(r\otimes\ell_{p_1}\wedge\cdots\wedge\ell_{p_n})
  =\sum_{i=1}^n(-1)^{i-1}rx_{p_i}\otimes
       \ell_{p_1}\wedge\cdots\widehat{\ell_{p_i}}\cdots
                              \wedge\ell_{p_n}.
\end{equation}
The indices here satisfy $p_1<\cdots<p_n$. The two orders in
which a pair of indices is removed have opposite signs and
the coefficients commute, so $d^2=0$.

This is a free resolution of the augmentation module $\mathbb C$.
For a finite set of primes the sequence of its $x_p$ is regular:
after any preceding variables have been set to one, the quotient
is still a Laurent polynomial algebra in the remaining variables,
and the next $t_p-1$ is a nonzero divisor. The finite Koszul complex
is exact in positive degrees by induction: adjoining one variable
is the mapping cone of multiplication by that variable's $x_p$,
whose kernel on the preceding degree-zero quotient is zero and
whose cokernel is the further quotient. The full complex ABR20 is
the filtered union of these finite complexes. A cycle uses finitely
many wedge indices; it is a boundary already in the corresponding
finite complex. Degree zero is $R/(x_p:p\text{ prime})=\mathbb C$.
This proves exactness also for the infinitely generated group, with
only finite algebraic sums at every step.

Tensor this resolution with the coefficient modules in the exact
sequence
\[
 0\longrightarrow M_0\longrightarrow R\otimes H
          \xrightarrow{\varepsilon\otimes m}\mathbb C^2
          \longrightarrow0.
\]
Here $\mathbb C^2$ has trivial $R$-action. Tensoring each free
Koszul degree preserves the short exact sequence. For the middle
module, its Koszul homology is $H$ in degree zero and zero in
positive degrees, because the complex is $K_\bullet(R)\otimes H$.
For the last module its differential is zero, and its degree-$n$
homology is $\Lambda^nW\otimes\mathbb C^2$. Consequently
\begin{equation}\tag{ABR21}
 H_n(G,M_0)\cong\Lambda^{n+1}W\otimes\mathbb C^2
            \quad(n\ge1),
 \qquad H_0(G,M_0)=Q_0,
\end{equation}
with ABR11 as the degree-zero exact sequence.

The exact connecting representative of an ordered element is
\begin{equation}\tag{ABR22}
 \begin{split}
 &\partial_n\bigl((\ell_{p_1}\wedge\cdots\wedge\ell_{p_{n+1}})
                                                    \otimes u\bigr)\\
 &\quad=\left[
  \sum_{i=1}^{n+1}(-1)^{i-1}
       (x_{p_i}\otimes\sigma(u))\otimes
       \ell_{p_1}\wedge\cdots\widehat{\ell_{p_i}}\cdots
                                  \wedge\ell_{p_{n+1}}
         \right]\in H_n(G,M_0).
 \end{split}
\end{equation}
Each coefficient is in $M_0$ because $\varepsilon(x_p)=0$.
This representative is the differential of
$(t_1\otimes\sigma(u))\otimes
  \ell_{p_1}\wedge\cdots\wedge\ell_{p_{n+1}}$
in the middle module's complex, so it is a cycle. Changing a
lift of $u$ changes it by a boundary in the first complex.
For $n\ge1$ its map is an isomorphism: a cycle of the first
complex becomes a boundary in the acyclic positive-degree middle
complex, proving surjectivity; if a connecting class is a boundary,
subtract its first-complex primitive from the chosen lift and use
middle-complex acyclicity one degree higher, proving injectivity.
Projection to the last complex has zero differential, so the latter
argument forces the original exterior class to be zero. For $n=0$
the representative is $[(t_p-1)\otimes\sigma(u)]$, exactly $j$ in
ABR11. These proofs retain the original wedge order and sign.

The module $V^K$ has trivial full-rational action, and its Koszul
homology is $\Lambda^nW\otimes V^K$. Periodization kills every
coefficient $x_p\otimes h$ in ABR22 separately, by ABR12. Thus
the homology map induced by $E\Theta:M_0\to V^K$ is zero in
every positive degree. In degree zero it is exactly ABR17.
The action $\mathscr R_b$ commutes with the differential ABR20.
Under the connecting isomorphism ABR21 its action is precisely
$\operatorname{id}_{\Lambda^{n+1}W}\otimes D_b$. To verify the
section-independence directly in ABR22, replace $R_b\sigma(u)$
by $\sigma(D_bu)$. Their difference is in $H_{00}$, and the
resulting difference of cycles is the differential of
$(t_1\otimes(R_b\sigma(u)-\sigma(D_bu)))$
times the complete ordered wedge, which lies in the first
complex itself. Thus every higher connecting class has the same
two original endpoint characters, with its prime wedge unchanged.

\subsection{The existing mass and endpoint sources under this map}

The actual source spaces, moments, group action, quotient and
coordinate maps in ABR1 and ABR6--10 agree with MRT1--5 and
ER27--30, not just their dimensions. The MRT source heat kernel is
\[
 k_T(x)=(4\pi T)^{-1/2}e^{-x^2/(4T)},\qquad T>0.
\]
For its even finite coefficient source $a$,
$A_Ta(x)=\sum_na_nk_T(x-n)$ is even Schwartz and has exactly
\begin{equation}\tag{ABR23}
 m(A_Ta)=\left(\sum_na_nk_T(n),\sum_na_n\right)
                  =(q_T(a),m(a)).
\end{equation}
Therefore the actual class and antecedent are
\begin{equation}\tag{ABR24}
 \begin{split}
 B_{p,T}(a)&=[(t_p-1)\otimes A_Ta],\\
 \Theta((t_p-1)\otimes A_Ta)(x)
  &=(A_Ta)(x_\infty/p)c_p(x_f)-(A_Ta)(x_\infty)c_1(x_f),\\
 (\rho,\beta)B_{p,T}(a)
  &=\left(0,\ell_p\otimes(q_T(a),m(a))\right).
 \end{split}
\end{equation}
Each such antecedent has both moments zero, including the separate
finite Haar factor $p^{-1}$ and the real integral factor $p$ of
its first summand. Its two periodizations are equal term by term
after the exact arithmetic reindexing ABR12.

These finite Gaussian sources attain all of $\mathcal B_{\mathrm{pr}}$.
For a prescribed pair $(u_0,u_1)$ and a fixed $T>0$, retain
$k_0=k_T(0)$ and $k_1=k_T(1)$ and take
\[
 a_0=\frac{u_0-k_1u_1}{k_0-k_1},\qquad
 a_1=a_{-1}=\frac{k_0u_1-u_0}{2(k_0-k_1)},\qquad
 a_n=0\quad(|n|>1).
\]
Here $k_0>k_1$ and direct substitution gives ABR23 equal to
$(u_0,u_1)$. Summing the finitely many selected prime classes
therefore realizes every boundary vector, with the original
kernel and coordinate conventions of MRT8.

For the Gaussian coefficient family $a_n^x=e^{-\pi n^2x}$,
$x>0$, MRT13--14 proves that the same $A_Ta^x$ is even Schwartz
and gives
\begin{equation}\tag{ABR25}
 \beta B_{p,T}(a^x)
 =\ell_p\otimes\left(
 (4\pi T)^{-1/2}\vartheta(x+1/(4\pi T)),\ \vartheta(x)\right),
 \qquad \rho B_{p,T}(a^x)=0.
\end{equation}
No heat time has been identified with the independent parameter $x$.
The separately specified MRT supported-zero source operation
$\widehat E$ still has its actual receiving matrix
$J_T=\left(\begin{smallmatrix}0&k_T(0)\\0&1\end{smallmatrix}\right)$
on these two moments, as proved in MRT11. Its periodization is
zero before and after that operation. The present comparison does
not redefine $\widehat E$ as the scalar supported-zero operation
on a different common-label vector carrier.

For an arbitrary completed MRT input $(a,0,m)$, MRT10 already
defines its boundary through $(q_T(a),m)$. ABR10 supplies the
Schwartz representative $\sigma(q_T(a),m)$ and hence an actual
adelic antecedent for this boundary. This does not assert that
$A_Ta$ itself is Schwartz for every completed coefficient sequence.

For the ER source, keep its original real variable $Z$, time $t>0$,
and function
\begin{equation}\tag{ABR26}
 h_t(Z)=e^{-Z^2/(4t)}H_t(Z)\cos(Z/(2t)).
\end{equation}
It is even Schwartz. The exact moments are
\begin{equation}\tag{ABR27}
 m(h_t)=\left(H_t(0),\frac{\sqrt{\pi t}}8e^{-1/(4t)}\right).
\end{equation}
To retain the reason for the second constant, insert the original
cosine integral $H_t(Z)=\int_0^\infty e^{tu^2}\Phi(u)\cos(Zu)\,du$.
The inner Gaussian integral is
\[
 \sqrt{\pi t}\{e^{-t(u+1/(2t))^2}+e^{-t(u-1/(2t))^2}\}
 =2\sqrt{\pi t}\,e^{-tu^2-1/(4t)}\cosh u.
\]
Absolute Fubini follows from the Gaussian and the faster-than-Gaussian
decay of the original $\Phi$. The retained factor $e^{tu^2}$
then gives $2\sqrt{\pi t}e^{-1/(4t)}H_0(i)$.
ER3--5 gives $H_0(-i)=1/16$ from the original limit
\[
 \frac1{16}\lim_{s\to1}
 s(s-1)\pi^{-s/2}\Gamma(s/2)\zeta(s)=\frac1{16};
\]
evenness gives $H_0(i)=H_0(-i)$. This proves ABR27 with the
original zeta pole, Gamma factor, power of pi and factor 16 retained.
No arithmetic Hurwitz-flow time is being substituted for $t$.

In ABR3 use this same real variable for $x_\infty$ and $h_t$.
The ER prime class has the actual antecedent and receiver
\begin{equation}\tag{ABR28}
 \begin{split}
 \Xi_{p,t}(x)
   &=h_t(x_\infty/p)c_p(x_f)-h_t(x_\infty)c_1(x_f),\\
 \rho[(t_p-1)\otimes h_t]&=0,\\
 \beta[(t_p-1)\otimes h_t]
   &=\ell_p\otimes
       \left(H_t(0),\frac{\sqrt{\pi t}}8e^{-1/(4t)}\right)\ne0.
 \end{split}
\end{equation}
The last strict nonzero statement uses $t>0$. Thus it is a
nonzero element of the exact kernel ABR17, with both original
source moments computed rather than assumed.

\subsection{Common-label support and the full SWQ receiving quotient}

For every complex vector space $U$ keep
\[
 G_L(U)=\{(0,\lambda):\lambda\in L\}\cup(U\times\{1_L\})
\]
for the original arbitrary nontrivial bounded distributive lattice.
Every linear map $T$ lifts by
$T^\uparrow(u,\lambda)=(Tu,\lambda)$, preserving addition with
label join, scalar action with label meet, and compositions.
The exact coinvariant quotient here is the same-label congruence
\begin{equation}\tag{ABR29}
 (m,\lambda)\sim(m',\lambda')
 \quad\Longleftrightarrow\quad
 \lambda=\lambda'\ \hbox{ and }\ m-m'\in IM_0.
\end{equation}
Its quotient is $G_L(Q_0)$. The further map $\rho^\uparrow$ has
fibres with the same labels and amplitude differences in
$j(\mathcal B_{\mathrm{pr}})$. Therefore a top-labelled nonzero
prime class from ABR24 or ABR28 maps to $(0,1_L)=e$, not
to $\tau=(0,0_L)$. For a general zero-amplitude label its image
is $(0,\lambda)$ with that same label. The joint map
$\Omega^\uparrow$ is an isomorphism onto
$G_L(V^K\oplus\mathcal B_{\mathrm{pr}})$; the entire direct sum
has one common label, not independently chosen component labels.

The Koszul amplitude complexes may be retained with these same
common-label lifts. Their composites satisfy
\begin{equation}\tag{ABR30}
 d^\uparrow d^\uparrow(z,\lambda)=(0,\lambda)
                                  =e\cdot(z,\lambda).
\end{equation}
Thus they are not ordinary complexes with a constant-$\tau$
zero map. The amplitude-zero cycles are exactly $G_L(\ker d)$,
and the same-label congruence by $\operatorname{im}d$ gives
$G_L(H_n(G,M_0))$. Under the induced periodization the lifted
classes ABR22 map to their retained zero labels in all positive
degrees. This states the receiving maps without discarding the
difference between amplitude zero and global absence.

Finally the range $V^K$ in ABR17 is contained in the actual
trivial-correspondence image $V$, not in a complementary spectral
space. Let $B_{\rm CC}$ denote the ambient strong test algebra
of SWQ and keep its original quotient
$Q_{\rm CC}=B_{\rm CC}/\overline V$. The further composition
\[
 Q_0\xrightarrow{\rho}V^K\longrightarrow B_{\rm CC}
                            \longrightarrow Q_{\rm CC}
\]
is the zero linear map because $V^K\subset V$. Its common-label
lift retains each $(0,\lambda)$, in accordance with SWQ's
same-label quotient; it does not identify $e$ and $\tau$.
The exact faithful source receiver is ABR18, which retains the
boundary rather than asserting an injection into this spectral
quotient. No assertion that $V$ is closed, no new metric on the
prime boundary, and no positivity or Riemann-hypothesis theorem
is made by this calculation.
